\scanpage{17}
\mspage{8}
Le fait d'avoir élargi le \uline{groupoïde} de
Teichmüller-Poincaré en la catégorie (qui n'est
plus un groupoïde !) de Teichmüller-
Galois précédente, est une façon \deleted{\unsure{commode}} \inserted{naturelle}
pour relier entre eux des courbes alg. relatifs
à des types numériques $(g,\nu)$ différents.
Une part essentielle de mon ``programme
à l'horizon'' est de déterminer de façon
plus ou moins complète cette catégorie
de Galois-Teichmüller, \inserted{par voie ``algébrique'' (ni transcendante ni de géométrie algébrique !)} \uline{avec} les \uline{opérations}
de $\operatorname{Aut}k$ \uline{dessus}, tout au moins dans
le cas où $k=\overline{\mathbb Q}$, et aussi (dans une
certaine mesure) lorsqu'on prend des
$k=\overline K$, $K$ corps de type fini sur $\unsure{\mathbb Q|Q}$, et
qu'on se limite aux \deleted{$\unsure{K}$}-automs \inserted{de $\overline K$} qui
sont des $K$-automs.\ldots\ \unsure{(?)|(!)}.
\margin{``algébrique'' \unsure{devant} signifier ici : par des méthodes de topologie \unsure{surfaces}, de \uline{combinatoire} (\unsure{courbes|cartes}) et de méthodes de théorie des groupes profinis\ldots}
\margin{16}
